\documentclass[UTF8,11pt,a4paper]{ctexart}
\usepackage[left=2.2cm,right=2.2cm,top=2.2cm,bottom=2.2cm]{geometry}
\usepackage{amsmath,amssymb,bm}
\usepackage{booktabs,longtable,array,multirow}
\usepackage{enumitem}
\usepackage{xcolor}
\usepackage{hyperref}
\usepackage{fancyhdr}
\usepackage{titlesec}
\usepackage{tikz}
\usetikzlibrary{arrows.meta,positioning,fit,shapes.geometric,calc}
\usepackage{tcolorbox}
\usepackage{microtype}
\usepackage{setspace}
\usepackage{caption}
\usepackage{listings}

\definecolor{DeepBlue}{RGB}{35,72,120}
\definecolor{LightBlue}{RGB}{235,243,252}
\definecolor{DeepGreen}{RGB}{44,110,82}
\definecolor{LightGreen}{RGB}{236,248,242}
\definecolor{WarmGray}{RGB}{245,245,245}
\definecolor{Accent}{RGB}{168,74,43}

\hypersetup{colorlinks=true,linkcolor=DeepBlue,urlcolor=DeepBlue,citecolor=DeepBlue}
\pagestyle{fancy}
\fancyhf{}
\lhead{DiT 在机器人动作生成中的完整讲解}
\rhead{架构、训练与推理}
\cfoot{\thepage}
\setlength{\headheight}{14pt}
\setstretch{1.18}
\setlist[itemize]{leftmargin=2em,itemsep=0.2em,topsep=0.3em}
\setlist[enumerate]{leftmargin=2em,itemsep=0.25em,topsep=0.3em}
\titleformat{\section}{\Large\bfseries\color{DeepBlue}}{\thesection}{0.8em}{}
\titleformat{\subsection}{\large\bfseries\color{DeepGreen}}{\thesubsection}{0.7em}{}

\newtcolorbox{keybox}[1][]{colback=LightBlue,colframe=DeepBlue,boxrule=0.7pt,arc=2mm,left=2mm,right=2mm,top=1.5mm,bottom=1.5mm,#1}
\newtcolorbox{examplebox}[1][]{colback=LightGreen,colframe=DeepGreen,boxrule=0.7pt,arc=2mm,left=2mm,right=2mm,top=1.5mm,bottom=1.5mm,#1}
\newtcolorbox{warningbox}[1][]{colback=WarmGray,colframe=Accent,boxrule=0.7pt,arc=2mm,left=2mm,right=2mm,top=1.5mm,bottom=1.5mm,#1}

\tikzset{
  block/.style={draw=DeepBlue,rounded corners=2mm,minimum width=3.1cm,minimum height=0.95cm,align=center,fill=LightBlue,font=\small},
  greenblock/.style={draw=DeepGreen,rounded corners=2mm,minimum width=3.1cm,minimum height=0.95cm,align=center,fill=LightGreen,font=\small},
  smallblock/.style={draw=black!55,rounded corners=1.5mm,minimum width=2.5cm,minimum height=0.75cm,align=center,fill=white,font=\footnotesize},
  arrow/.style={-{Latex[length=2.4mm]},thick,draw=black!70}
}

\title{\textbf{DiT 在机器人动作生成中的完整讲解}\\[0.4em]
\large 从完整架构到训练流程与推理流程}
\author{技术讲解报告}
\date{\today}

\begin{document}
\maketitle

\begin{keybox}
\textbf{本文目标：}以 GR00T N1 一类视觉-语言-动作系统为背景，系统解释用于动作生成的 DiT（Diffusion Transformer）。全文严格按照三个问题展开：
\begin{enumerate}
  \item 完整的 DiT 架构包含哪些模块，各模块输入、输出和作用是什么；
  \item 完整训练流程如何进行；
  \item 完整推理流程如何进行。
\end{enumerate}
\end{keybox}

\tableofcontents
\newpage

\section{整体认识：DiT 在 VLA 系统中的位置}

在视觉-语言-动作（Vision-Language-Action, VLA）系统中，可以把系统粗略分成两个层次：

\begin{center}
\begin{tikzpicture}[node distance=1.15cm]
\node[block] (input) {图像 + 语言指令};
\node[block,below=of input] (vlm) {VLM：理解场景与任务};
\node[greenblock,below=of vlm] (latent) {视觉语言 Token / Latent};
\node[block,below=of latent] (dit) {DiT：生成连续动作};
\node[greenblock,below=of dit] (action) {未来 $H$ 步机器人动作};
\draw[arrow] (input)--(vlm);
\draw[arrow] (vlm)--(latent);
\draw[arrow] (latent)--(dit);
\draw[arrow] (dit)--(action);
\end{tikzpicture}
\end{center}

VLM 主要负责回答：看到了什么、目标是什么、指令要求做什么、物体之间有哪些空间关系。DiT 则负责回答：机器人在当前状态下，身体具体应该怎样连续运动。

\begin{keybox}
\textbf{一句话概括：}VLM 决定“要做什么”，DiT 决定“身体具体怎么动”。
\end{keybox}

DiT 全称为 \textbf{Diffusion Transformer}。它并不等于“扩散过程”本身，而是扩散模型或 Flow Matching 中负责预测去噪方向的 Transformer 网络：

\[
\text{扩散/Flow Matching 规则} + \text{Transformer 网络} = \text{DiT 生成模型}
\]

其中：
\begin{itemize}
  \item 扩散或 Flow Matching 规定如何构造带噪样本、模型应预测什么、损失怎样计算、推理时如何从噪声生成数据；
  \item Transformer 负责根据当前带噪动作和条件信息，预测下一步应如何修正动作。
\end{itemize}

\section{符号与张量形状}

\subsection{机器人当前状态}
记当前机器人状态为
\[
s_t\in\mathbb{R}^{B\times d_s},
\]
其中 $B$ 是 Batch Size，$d_s$ 是状态维度。状态可以包括关节角、关节速度、末端位姿、夹爪开合程度等。

\subsection{动作块}
模型一次预测未来 $H$ 步动作：
\[
A_t=[a_t,a_{t+1},\ldots,a_{t+H-1}],
\]
其形状为
\[
A_t\in\mathbb{R}^{B\times H\times d_a}.
\]
这里：
\begin{itemize}
  \item $B$：一次处理多少条样本；
  \item $H$：动作块长度或预测时域；
  \item $d_a$：单步动作维度。
\end{itemize}

\begin{examplebox}
若 $B=32$、$H=16$、$d_a=8$，则张量形状为 $[32,16,8]$，表示一次处理 32 条样本，每条样本含未来 16 步动作，每一步动作是 8 维。
\end{examplebox}

\subsection{VLM 条件 Token}
VLM 输出一串视觉语言 Token：
\[
C\in\mathbb{R}^{B\times N_c\times d_c},
\]
其中 $N_c$ 是 Token 数量，$d_c$ 是 VLM 隐藏维度。这里使用的是 VLM 中间层隐藏表示，而不是已经经过 LM Head 的文字概率。

\subsection{Flow 时间}
记 Flow 时间为
\[
\tau\in[0,1].
\]
它表示当前动作样本处于“更接近纯噪声”还是“更接近真实动作”的阶段。

\section{完整的 DiT 架构}

一个用于机器人动作生成的完整 DiT 可以拆为六个核心模块：
\begin{enumerate}
  \item Action Encoder；
  \item State Encoder；
  \item Time Encoder；
  \item VLM 条件投影模块；
  \item DiT Transformer Backbone；
  \item Action Decoder / Output Head。
\end{enumerate}

\subsection{整体结构图}

\begin{center}
\resizebox{0.96\textwidth}{!}{%
\begin{tikzpicture}[node distance=0.85cm and 1.2cm]
\node[smallblock] (noisy) {带噪动作 $A_\tau$};
\node[smallblock,below=of noisy] (aenc) {Action Encoder};
\node[smallblock,below=of aenc] (atok) {动作 Tokens};

\node[smallblock,right=2.0cm of noisy] (state) {当前状态 $s_t$};
\node[smallblock,below=of state] (senc) {State Encoder};
\node[smallblock,below=of senc] (semb) {State Embedding};

\node[smallblock,right=2.0cm of state] (time) {Flow 时间 $\tau$};
\node[smallblock,below=of time] (tenc) {Time Encoder};
\node[smallblock,below=of tenc] (temb) {Time Embedding};

\node[smallblock,right=2.0cm of time] (cond) {VLM Tokens $C$};
\node[smallblock,below=of cond] (proj) {Condition Projection};
\node[smallblock,below=of proj] (ctok) {条件 Tokens};

\node[block,below=2.0cm of $(atok)!0.5!(ctok)$] (backbone) {DiT Transformer Blocks $\times N$\\Self-Attention + Cross-Attention + MLP + 条件调制};
\node[block,below=of backbone] (head) {Action Decoder / Output Head};
\node[greenblock,below=of head] (vel) {预测速度场 $v_\theta$};

\draw[arrow] (noisy)--(aenc); \draw[arrow] (aenc)--(atok);
\draw[arrow] (state)--(senc); \draw[arrow] (senc)--(semb);
\draw[arrow] (time)--(tenc); \draw[arrow] (tenc)--(temb);
\draw[arrow] (cond)--(proj); \draw[arrow] (proj)--(ctok);
\draw[arrow] (atok)--(backbone);
\draw[arrow] (semb)--(backbone);
\draw[arrow] (temb)--(backbone);
\draw[arrow] (ctok)--(backbone);
\draw[arrow] (backbone)--(head);
\draw[arrow] (head)--(vel);
\end{tikzpicture}}
\end{center}

\subsection{模块一：Action Encoder}

\textbf{输入：}
\[
A_\tau\in\mathbb{R}^{B\times H\times d_a}.
\]
它是当前时刻的带噪动作块。

\textbf{作用：}把每一步低维机器人动作映射到 Transformer 的统一隐藏维度：
\[
\mathbb{R}^{d_a}\rightarrow\mathbb{R}^{D}.
\]

\textbf{输出：}
\[
X_a\in\mathbb{R}^{B\times H\times D}.
\]
每一个未来动作都变成一个动作 Token。

\begin{examplebox}
当前例子中，若机器人每步动作为 7 个关节控制量加 1 个夹爪控制量，则 $d_a=8$。Action Encoder 可将 $[B,16,8]$ 映射为 $[B,16,1024]$。
\end{examplebox}

\subsection{模块二：State Encoder}

\textbf{输入：}
\[
s_t\in\mathbb{R}^{B\times d_s}.
\]

\textbf{作用：}告诉 DiT 机器人当前的物理状态。相同任务在不同起始关节角或不同末端位置下，需要不同动作轨迹。

\textbf{输出：}
\[
e_s\in\mathbb{R}^{B\times D}.
\]
状态表示可作为状态 Token、与动作 Token 拼接、通过 Cross-Attention 注入，或用于条件归一化调制。

\subsection{模块三：Time Encoder}

\textbf{输入：}$\tau\in[0,1]$。

\textbf{作用：}告诉模型当前噪声阶段。相同数值在高噪声阶段可能需要大幅修正，在低噪声阶段则只需轻微修正。

典型结构为：
\[
\tau\rightarrow\text{Fourier/Sinusoidal Encoding}\rightarrow\text{MLP}\rightarrow e_\tau.
\]

\textbf{输出：}
\[
e_\tau\in\mathbb{R}^{B\times D}.
\]

时间表示通常通过 AdaLN（Adaptive LayerNorm）、加法、门控或额外条件 Token 注入每个 DiT Block。

\subsection{模块四：VLM 条件投影}

\textbf{输入：}
\[
C\in\mathbb{R}^{B\times N_c\times d_c}.
\]

\textbf{作用：}把 VLM Token 映射到与 DiT 兼容的维度：
\[
\mathbb{R}^{d_c}\rightarrow\mathbb{R}^{D}.
\]

\textbf{输出：}
\[
C'\in\mathbb{R}^{B\times N_c\times D}.
\]

随后动作 Token 通过 Cross-Attention 读取这些视觉语言条件：
\[
\operatorname{Attention}(Q_{\text{action}},K_{\text{VLM}},V_{\text{VLM}}).
\]

其中 Query 来自动作 Token，Key 和 Value 来自 VLM Token。它让每一步动作都能关注目标物体、指令语义和空间关系。

\subsection{模块五：DiT Transformer Backbone}

Backbone 由 $N$ 个 DiT Block 堆叠构成。一个简化 Block 包含：

\begin{center}
\begin{tikzpicture}[node distance=0.72cm]
\node[smallblock] (x) {输入动作 Tokens};
\node[smallblock,below=of x] (ln1) {LayerNorm / AdaLN};
\node[smallblock,below=of ln1] (sa) {Self-Attention};
\node[smallblock,below=of sa] (res1) {残差连接};
\node[smallblock,below=of res1] (ln2) {LayerNorm / AdaLN};
\node[smallblock,below=of ln2] (ca) {Cross-Attention};
\node[smallblock,right=1.5cm of ca] (vlmc) {VLM 条件 Tokens};
\node[smallblock,below=of ca] (res2) {残差连接};
\node[smallblock,below=of res2] (ln3) {LayerNorm / AdaLN};
\node[smallblock,below=of ln3] (ffn) {MLP / FFN};
\node[smallblock,below=of ffn] (out) {输出动作 Tokens};
\foreach \a/\b in {x/ln1,ln1/sa,sa/res1,res1/ln2,ln2/ca,ca/res2,res2/ln3,ln3/ffn,ffn/out}{\draw[arrow] (\a)--(\b);}
\draw[arrow] (vlmc)--(ca);
\end{tikzpicture}
\end{center}

\paragraph{Self-Attention}建模动作序列内部的时间依赖。例如第 12 步动作必须知道前 11 步已把机械臂移动到哪里，才能保持轨迹连续并在正确时刻闭合夹爪。

\paragraph{Cross-Attention}让动作 Token 读取 VLM 条件，例如“红色杯子在哪里”“要抓取还是放置”“目标容器在哪一侧”。

\paragraph{MLP / FFN}对每个 Token 进行非线性变换，典型维度为 $D\rightarrow 4D\rightarrow D$。

\paragraph{残差连接}使深层 Transformer 更容易优化：
\[
x_{l+1}=x_l+F(x_l).
\]

\subsection{模块六：Action Decoder / Output Head}

\textbf{输入：}
\[
X_{\text{out}}\in\mathbb{R}^{B\times H\times D}.
\]

\textbf{作用：}将 Transformer 隐藏维度映射回动作维度：
\[
\mathbb{R}^{D}\rightarrow\mathbb{R}^{d_a}.
\]

\textbf{输出：}
\[
\hat V\in\mathbb{R}^{B\times H\times d_a}.
\]

在 Flow Matching 中，这通常不是最终动作，而是速度场：
\[
v_\theta(A_\tau,\tau,s_t,C).
\]
它表示当前动作样本应该向哪个方向移动，以及移动多快。

\section{完整训练流程}

一条训练样本写作：
\[
(I_t,l,s_t,A_1),
\]
其中 $I_t$ 是图像，$l$ 是语言指令，$s_t$ 是当前状态，$A_1$ 是真实未来动作块。

\subsection{步骤 1：读取专家轨迹样本}
例如：
\begin{itemize}
  \item 图像：桌面上存在一个红色杯子；
  \item 指令：Pick up the red cup；
  \item 状态：当前关节角、速度和夹爪状态；
  \item 动作：未来 16 步专家动作。
\end{itemize}

真实动作形状为
\[
A_1\in\mathbb{R}^{B\times H\times d_a}.
\]

\subsection{步骤 2：VLM 提取视觉语言条件}
图像与语言经过 VLM，取中间层视觉语言 Token：
\[
C=\operatorname{VLM}(I_t,l).
\]

\subsection{步骤 3：编码当前状态}
\[
e_s=\operatorname{StateEncoder}(s_t).
\]

\subsection{步骤 4：采样随机噪声}
采样与真实动作形状完全相同的高斯噪声：
\[
A_0\sim\mathcal N(0,I),\qquad A_0\in\mathbb{R}^{B\times H\times d_a}.
\]

\subsection{步骤 5：采样 Flow 时间}
\[
\tau\sim U(0,1).
\]

\subsection{步骤 6：构造中间带噪动作}
在线性 Flow Matching 的简单形式中：
\[
A_\tau=(1-\tau)A_0+\tau A_1.
\]
当 $\tau=0$ 时，$A_\tau=A_0$，即纯噪声；当 $\tau=1$ 时，$A_\tau=A_1$，即真实动作。

\subsection{步骤 7：构造监督速度}
对上式关于 $\tau$ 求导：
\[
\frac{dA_\tau}{d\tau}=A_1-A_0.
\]
因此真实速度目标为：
\[
v^*=A_1-A_0.
\]

\subsection{步骤 8：DiT 前向传播}
模型接收四类输入：
\[
(A_\tau,\tau,s_t,C),
\]
输出预测速度：
\[
\hat v=v_\theta(A_\tau,\tau,s_t,C).
\]

\subsection{步骤 9：计算训练损失}
最常见的是均方误差：
\[
\mathcal L_{\mathrm{FM}}=
\left\|v_\theta(A_\tau,\tau,s_t,C)-(A_1-A_0)\right\|_2^2.
\]

\subsection{步骤 10：反向传播与参数更新}
根据损失反向传播，更新 Action Encoder、State Encoder、Time Encoder、条件投影、DiT Backbone 和 Action Decoder。若采用端到端联合训练，也可更新 VLM 的部分参数。

\subsection{训练流程总图}
\begin{center}
\resizebox{0.9\textwidth}{!}{%
\begin{tikzpicture}[node distance=0.8cm and 1.3cm]
\node[block] (data) {训练样本：图像、语言、状态、真实动作 $A_1$};
\node[smallblock,below left=1.0cm and 1.7cm of data] (vlm) {VLM 提取条件 $C$};
\node[smallblock,below=of data] (noise) {采样噪声 $A_0$};
\node[smallblock,below right=1.0cm and 1.7cm of data] (tau) {采样时间 $\tau$};
\node[smallblock,below=1.0cm of noise] (interp) {$A_\tau=(1-\tau)A_0+\tau A_1$};
\node[block,below=1.2cm of interp] (dit) {DiT：输入 $A_\tau,\tau,s_t,C$};
\node[greenblock,below=of dit] (pred) {预测速度 $v_\theta$};
\node[smallblock,below=of pred] (loss) {与 $A_1-A_0$ 计算 MSE};
\node[smallblock,below=of loss] (bp) {反向传播与更新参数};
\draw[arrow] (data)--(vlm);
\draw[arrow] (data)--(noise);
\draw[arrow] (data)--(tau);
\draw[arrow] (noise)--(interp);
\draw[arrow] (tau)--(interp);
\draw[arrow] (vlm)--(dit);
\draw[arrow] (interp)--(dit);
\draw[arrow] (dit)--(pred);
\draw[arrow] (pred)--(loss);
\draw[arrow] (loss)--(bp);
\end{tikzpicture}}
\end{center}

\section{完整推理流程}

推理阶段与训练阶段的根本区别是：\textbf{没有真实未来动作 $A_1$}。因此模型必须从随机动作噪声开始生成。

\subsection{步骤 1：获取当前观测}
机器人读取：
\[
I_t,\quad l,\quad s_t.
\]
即当前图像、语言指令与机器人状态。

\subsection{步骤 2：VLM 提取条件}
\[
C=\operatorname{VLM}(I_t,l).
\]

\subsection{步骤 3：初始化随机动作块}
\[
A^{(0)}\sim\mathcal N(0,I).
\]
例如单条样本下：
\[
A^{(0)}\in\mathbb{R}^{1\times 16\times 8}.
\]
此时它是纯随机动作，不能直接执行。

\subsection{步骤 4：设置去噪/积分时间网格}
将 $[0,1]$ 划分为 $K$ 个步骤：
\[
0=\tau_0<\tau_1<\cdots<\tau_K=1.
\]

\subsection{步骤 5：反复调用 DiT 预测速度}
在第 $k$ 步：
\[
v^{(k)}=v_\theta(A^{(k)},\tau_k,s_t,C).
\]

\subsection{步骤 6：数值积分更新动作}
最简单的 Euler 更新为：
\[
A^{(k+1)}=A^{(k)}+\Delta\tau_k\,v^{(k)},
\]
其中
\[
\Delta\tau_k=\tau_{k+1}-\tau_k.
\]
重复 $K$ 次后，随机动作逐步变成合理动作。

\subsection{步骤 7：得到最终动作块}
\[
\hat A_t=A^{(K)}=[\hat a_t,\hat a_{t+1},\ldots,\hat a_{t+H-1}].
\]

\subsection{步骤 8：执行与重新规划}
常见执行方式有三种：
\begin{enumerate}
  \item 只执行第一步，然后重新感知并重新规划；
  \item 执行前若干步，再重新规划；
  \item 对连续多次预测做 Temporal Ensembling，以获得更平滑控制。
\end{enumerate}

\begin{warningbox}
“预测未来 16 步”并不意味着机器人一定开环执行全部 16 步。实践中通常采用滚动时域控制：预测一段、执行一小段、重新观测、再次预测。
\end{warningbox}

\subsection{推理流程总图}
\begin{center}
\resizebox{0.88\textwidth}{!}{%
\begin{tikzpicture}[node distance=0.78cm]
\node[block] (obs) {当前图像 + 语言 + 状态};
\node[smallblock,below=of obs] (cond) {VLM / State Encoder 提取条件};
\node[smallblock,below=of cond] (init) {初始化随机动作 $A^{(0)}$};
\node[block,below=of init] (loop) {第 $k$ 次 DiT：预测 $v^{(k)}$};
\node[smallblock,below=of loop] (update) {$A^{(k+1)}=A^{(k)}+\Delta\tau_k v^{(k)}$};
\node[smallblock,below=of update] (repeat) {重复 $K$ 次};
\node[greenblock,below=of repeat] (chunk) {得到未来 $H$ 步动作};
\node[smallblock,below=of chunk] (exec) {执行第一步或前几步};
\node[smallblock,below=of exec] (replan) {重新观测并再次规划};
\foreach \a/\b in {obs/cond,cond/init,init/loop,loop/update,update/repeat,repeat/chunk,chunk/exec,exec/replan}{\draw[arrow] (\a)--(\b);}
\draw[arrow] (repeat.west) .. controls +(-2.6,0) and +(-2.6,0) .. node[left,font=\footnotesize,align=center]{未到\\$K$ 次} (loop.west);
\end{tikzpicture}}
\end{center}

\section{训练与推理的核心区别}

\begin{longtable}{>{\raggedright\arraybackslash}p{0.18\textwidth}>{\raggedright\arraybackslash}p{0.36\textwidth}>{\raggedright\arraybackslash}p{0.36\textwidth}}
\toprule
对比项 & 训练阶段 & 推理阶段\\
\midrule
是否有真实动作 & 有真实未来动作 $A_1$ & 没有真实未来动作\\
初始动作 & 由噪声 $A_0$ 与真实动作 $A_1$ 插值得到 $A_\tau$ & 直接从随机高斯噪声 $A^{(0)}$ 开始\\
模型目标 & 预测已知的真实速度 $A_1-A_0$ & 预测当前速度并逐步更新动作\\
损失计算 & 使用 MSE 等监督损失 & 不计算训练损失\\
调用次数 & 每个训练样本通常随机取一个 $\tau$ 做一次前向传播 & 一次动作生成需要重复调用 DiT 多次\\
最终用途 & 更新网络参数 & 得到可执行动作块\\
\bottomrule
\end{longtable}

\section{为什么 DiT 能把噪声变成动作}

DiT 并不是天然具有“去噪能力”。它在训练中反复看到如下对应关系：

\begin{center}
\begin{tikzpicture}[node distance=0.9cm]
\node[block] (input) {带噪动作 + 时间 + 状态 + VLM 条件};
\node[greenblock,below=of input] (target) {正确速度方向 $A_1-A_0$};
\draw[arrow] (input)--(target);
\end{tikzpicture}
\end{center}

经过大量训练，模型学习到一个条件速度场：
\[
v_\theta(A,\tau,s,C).
\]
这个速度场可以理解为：在动作空间的任意位置，给定当前机器人状态和任务条件，告诉当前动作样本下一步应该向哪里移动。推理时从随机点出发，只要不断沿速度场前进，就会逐渐靠近真实动作分布。

\section{为什么不直接用 MLP 回归动作}

可以直接训练：
\[
\hat A=f(I,l,s),
\]
但直接使用 MSE 回归时，若同一状态存在多种合理轨迹，模型容易将它们平均。例如抓杯子可以从左侧接近，也可以从右侧接近，平均后可能得到一条并不合理的中间轨迹。

Flow Matching 更接近学习条件动作分布：
\[
p(A\mid I,l,s),
\]
因此更适合表达多种合理动作、连续动作序列以及多模态轨迹分布。

\section{与 LLM 语言输出的统一理解}

LLM 的输出过程是：
\[
\text{Hidden State}\rightarrow\text{LM Head}\rightarrow\text{词表 Logits}\rightarrow\text{Token}.
\]

DiT 的输出过程是：
\[
\text{动作 Hidden Tokens}\rightarrow\text{Action Head}\rightarrow\text{速度场}\rightarrow\text{数值积分}\rightarrow\text{动作序列}.
\]

二者都使用 Transformer 和输出 Head，但任务空间不同：LM Head 映射到词表，Action Head 映射到动作空间；LLM 自回归生成文字 Token，DiT 从噪声迭代生成动作。

\section{最终总结}

\begin{keybox}
\textbf{完整总结：}机器人动作 DiT 接收带噪动作块、Flow 时间、当前机器人状态和 VLM 视觉语言 Token，通过多层 Self-Attention、Cross-Attention、MLP 与条件调制，预测当前动作样本朝真实动作分布移动的速度方向。训练时利用真实动作和随机噪声构造监督速度场；推理时从高斯噪声出发，反复调用 DiT 并进行数值积分，最终生成未来一段连续机器人动作。机器人通常只执行其中第一步或前几步，再重新感知和规划，从而形成闭环控制。
\end{keybox}

\end{document}
